% Source: sources/returns/S6_04_MAPS_MANIFOLDS/sections/manifolds.tex:8957–9010
\paragraph{Theorem D.3.2 (Known, with derivation: the smooth side is $S^3\times S^3$).}
For the sphere twist data and $0<|s|\ll1$, the complex-oriented
smooth manifold $Y_{u,s}$ is orientation-preservingly diffeomorphic
to $S^3\times S^3$. Its integral homology is
$(\mathbb Z,0,0,\mathbb Z^2,0,0,\mathbb Z)$,
its middle intersection form is the rank-two unimodular
skew-symmetric form, and $w_2=p_1=0$.
\begin{proof}
An ordinary-double-point smoothing replaces a disc neighbourhood
of its exceptional curve by a disc neighbourhood of the vanishing
three-sphere. In the displayed coordinates this can be read
without a deformation-theoretic inference: the resolution is the
disc bundle of $\mathcal O(-1)^{\oplus2}$ over $S^2$, and the
Milnor fibre of the nondegenerate quadratic equation is the disc
cotangent bundle of $S^3$. The underlying oriented rank-four
bundle over $S^2$ is trivial, since its $w_2$ is the reduction of
$-2$ and oriented rank-four bundles over $S^2$ are classified by
$\pi_1SO(4)=\mathbb Z/2$. The latter disc cotangent bundle is
$D^3\times S^3$ because $S^3$ is parallelizable. Both boundaries
are $S^2\times S^3$.

Let $A$ be the complement of the interior of that resolution tube
in $X^{\log}(u)\simeq S^6$. General position makes $A$ simply connected:
a loop and a filling disc in $S^6$ can be kept disjoint from the
embedded two-sphere since $2+2<6$. Alexander duality gives
$\widetilde H_i(A;\mathbb Z)=\mathbb Z$ for $i=3$ and zero for
all other $i$. The normal three-sphere in the boundary generates
that group. In the Mayer--Vietoris sequence for
$Y=A\cup(D^3\times S^3)$, the degree-three map from the boundary
is therefore $\mathbb Z\to\mathbb Z^2$, with both entries units,
and the degree-two boundary group is $\mathbb Z$ with zero image.
It follows that $H_2(Y)=H_4(Y)=0$ and that $H_3(Y)$ is an extension
of $\mathbb Z$ by $\mathbb Z$, hence $\mathbb Z^2$.
Van Kampen gives $\pi_1(Y)=1$. The remaining homology follows by
duality, and the assertions about $w_2,p_1$ follow from the
vanishing of the corresponding cohomology groups.

We apply Wall's classification in the precise form of
Burklund--Hahn--Senger, Theorem 2.6, with their Definition 2.5 and
Recollection 2.2 \cite{BHS19}. This theorem says that two
$(n-1)$-connected closed $2n$-manifolds with the same middle
intersection form and normal invariant
$\alpha:H_n\to\pi_{n-1}SO(n)$ differ by a homotopy sphere.
Here $n=3$, the Hurewicz theorem gives $\pi_2(Y)=0$, and
$\pi_2SO(3)=0$, so $\alpha$ is zero. Every unimodular
skew-symmetric rank-two form has an oriented symplectic basis;
thus these invariants equal those of $S^3\times S^3$.
Finally $\Theta_6/bP_7=0$ by the dimension-six entry of
\cite[p.~512]{manKM63}, and $bP_7=0$ by its Theorem~5.1
(equivalently the framed-surgery Theorem~6.6). Thus $\Theta_6=0$
removes the homotopy-sphere ambiguity.
This derives the diffeomorphism assertion from the classification
rather than from Euler characteristic alone.
\end{proof}
